\begin{table*}[!ht]
\centering
\begin{tabular*}{\textwidth}{@{\extracolsep{\fill}}l >{\raggedright\arraybackslash}p{172pt} >{\raggedright\arraybackslash}p{172pt}@{}}
\toprule
\tabheadfont Item & \tabheadfont Single-nucleus RNA-seq & \tabheadfont Spatial transcriptomics \\
\midrule
Accession & GSE308103 & GSE307534 \\
Platform & fixed RNA profiling & Visium CytAssist, 11\,mm \\
Feature space & 18{,}069 genes analysed & 18{,}085 probes \\
Samples & 75 & 56 sections \\
Patients & 25 & 25 \\
\addlinespace[3pt]
Units called & 798{,}100 nuclei & 14{,}336 spots per section \\
Retained after QC & 767{,}839 (96.2\%) & 95.03\% of spots \\
Doublet rate & 15.48\% & --- \\
Analysed & \textbf{413{,}697 nuclei} & all 56 sections \\
Grid and spacing & --- & $128\times112$; 100\,$\mu$m \\
\addlinespace[3pt]
Patients in both modalities & \multicolumn{2}{l}{23} \\
\bottomrule
\multicolumn{3}{@{}l}{\footnotesize\itshape Sections per stage: Normal 1 \quad AAH 11 \quad AIS 14 \quad MIA 4 \quad IAC 26 (source token \texttt{LUAD})} \\
\end{tabular*}
\caption{\textbf{Cohort and data sources.} Both modalities are public, come from the same
patient series, and share 23 patients. Both are probe-based assays rather than
whole-transcriptome chemistries, so each has a defined feature space. The invasive stage is
labelled \texttt{LUAD} in the source annotations and is mapped to IAC throughout. Normal is
represented by a single spatial section and MIA by four, so estimates for these two stages
should not be quoted on their own.}
\label{tab:cohort}
\end{table*}
